% Source-recovered from the 2026-11-01 central-quotient manuscript, Sections 17--18.
% Input by papers/central_quotient_local_models.tex.
\section{中心値零の小次元表現と六次分裂元}
\label{sec:central-quotient-representations}

前節の代数$A=R_1$を
\[
 A=\CC\langle X,Y,Z\rangle/
 \bigl([X,Y]-Z^2,\ [Y,Z]-(X^2+1),\ [Z,X]-Y^2\bigr)
\]
とする。$\Omega_1=(X^3+Y^3+Z^3)/3
-(XYZ+XZY+YXZ+YZX+ZXY+ZYX)/9$、
$G_1=\Omega_1+X$、$g=3G_1/2$と置く。
本節の命題は旧論文VIの第17--18節から復元した。
関係式・$G_1$の明示表現上の消滅は、
\texttt{computations/verify\_central\_quotient\_reps.py}
により$\CC$の標数零の係数で厳密に検査できる。

\begin{proposition}[一次・二次元単純表現の不存在]
\label{prop:local-no-small-simples}
$g=0$上に一次元単純表現は存在せず、
$A$は二次元単純複素表現を持たない。
\end{proposition}
\begin{proof}
一次元表現では交換子は零なので
$Z=Y=0$、$X^2=-1$となる。
この場合$G_1=X^3/3+X=2X/3$で、
$g=X=\pm i\ne0$である。

二次元表現においてPauli行列$\boldsymbol{\sigma}$を使い
\[
 X=aI+\boldsymbol x\cdot\boldsymbol{\sigma},\quad
 Y=bI+\boldsymbol y\cdot\boldsymbol{\sigma},\quad
 Z=cI+\boldsymbol z\cdot\boldsymbol{\sigma}
\]
と書く。$\boldsymbol{x},\boldsymbol{y},\boldsymbol{z}\in\CC^3$とし
内積・外積は標準双線形形式を使う。
三つの関係式は
\[
 \begin{gathered}
 \boldsymbol x^2=-a^2-1,\quad
 \boldsymbol y^2=-b^2,\quad
 \boldsymbol z^2=-c^2,\\
 \boldsymbol x\times\boldsymbol y=-ic\boldsymbol z,\quad
 \boldsymbol y\times\boldsymbol z=-ia\boldsymbol x,\quad
 \boldsymbol z\times\boldsymbol x=-ib\boldsymbol y
 \end{gathered}
\]
と同値である。

$abc\ne0$なら三つのベクトルは互いに直交し、
$A_0=a^2,B_0=b^2,C_0=c^2$について
\[
 B_0(A_0+1)=C_0^2,\quad
 B_0C_0=A_0(A_0+1),\quad
 C_0(A_0+1)=B_0^2.
\]
第一・第三式から$(A_0+1)^2=B_0C_0$で、
第二式との比較は$A_0+1=A_0$を強制して矛盾する。
$abc=0$の場合は非零である$a,b,c$の個数で分ける。
二つ非零なら外積零と内積零により非零平方長のベクトルと
従属な零直交ベクトルから矛盾する。
$b$または$c$だけが非零の場合も
$\boldsymbol x^2=-1$と外積・内積の式から矛盾する。
$a$だけが非零の場合、$\boldsymbol x\ne0$なら
$\boldsymbol y,\boldsymbol z$はこれに比例し
$\boldsymbol y\times\boldsymbol z=-ia\boldsymbol x$に矛盾する。
$\boldsymbol x=0$なら$X=\pm iI$で、
$Y,Z$は互いに可換な平方零行列となり共通不変直線を持つ。
最後に$a=b=c=0$なら$X^2=-I$で
$[X,Y]=[Z,X]=0$だから$Y,Z$は$X$の固有基底で対角かつ
平方零、すなわち零である。この場合も可約。
従って二次元単純表現は存在しない。
\end{proof}

\subsection{軸上の二つの三次元単純表現}
\begin{proposition}[六次多項式による明示表現]
\label{prop:local-two-axial}
\[
 q(t)=7t^6+21t^4+14t^2+1,\quad
 \lambda_0=\frac{42t^5+119t^3+75t}{13},\quad
 \lambda_1=\frac{7t^5+35t^3+45t}{13},\quad
 \lambda_2=t
\]
とし、添字を$\mathbb Z/3$で読んで
$d_i=\lambda_{i+1}-\lambda_i$とする。
$q(t)=0$である任意の複素数$t$に対し
\[
 Xe_i=\lambda_i e_i,\qquad
 Ye_i=d_{i-1}e_{i+1},\qquad
 Ze_i=d_ie_{i-1}\qquad(i\in\mathbb Z/3)
\]
は$A/(G_1)$の三次元単純表現を与える。
$t$の六根から得られる同型類はちょうど二つであり、
$\operatorname{Tr}X=\pm\sqrt{-7}$で区別される。
両方の中心指標は$(g,U,V,W)=(0,1/7,0,0)$である。
両方において$Z^3$はスカラーであり、
二つの$Z$の固有値を合わせると$T^6+1/7$の六根となる。
\end{proposition}
\begin{proof}
上の行列に代入すれば三つの交換子関係と
$\Omega_1+X=0$が$\CC[t]/(q)$で成立する。
付属の厳密記号計算はすべての九行列成分の剰余零を検証する。
また
\[
 s=\lambda_0+\lambda_1+\lambda_2,\qquad
 s^2=-7,\quad
 d_0d_1d_2=-s/7,\quad
 Z^3=(-s/7)I
\]
が$q$での多項式剰余として成立する。
$q$は重根を持たず、$q$と$d_0d_1d_2$は互いに素である。
従って各$d_i\ne0$で、$X$は相異なる三固有値を持ち、
$Y$の三周期による推移性から表現は単純である。
$X$の三固有値は、$s$を固定したとき
$T^3-sT^2-2T+s/7$の三根である。
循環順序の違いは基底の循環置換で同一視される。
$s$は$\pm\sqrt{-7}$の二通りで、traceが異なる表現は同型でない。
二つでの$Z^3$はそれぞれ$\mp\sqrt{-7}/7$なので
固有値の六つの組は$T^6+1/7$の全ての根になる。
\end{proof}

\subsection{非軸上の三つの中心点を支配する行列}
原始三乗根$\omega^2+\omega+1=0$を固定し、
\[
 X_0=
 \begin{pmatrix}
 -1&1&1-\omega\\
 -1&(2\omega+1)/3&1\\
 -(\omega+2)/3&1/3&-(2\omega+1)/3
 \end{pmatrix},
 \quad Y_0=\operatorname{diag}(-1,\omega,\omega^2),
\]
\[
 Z_0=
 \begin{pmatrix}
 -1&-(2\omega+1)/3&\omega+1\\
 -\omega-1&(\omega-1)/3&\omega+1\\
 (\omega-1)/3&-\omega/3&-(\omega+2)/3
 \end{pmatrix}
\]
と置く。
\begin{proposition}[六つの非軸上表現]
\label{prop:local-six-nonaxial}
$\varepsilon=\pm1$、$k\in\mathbb Z/3$について
\[
 (X,Y,Z)=(\varepsilon X_0,\varepsilon\omega^kY_0,
                   \varepsilon\omega^{2k}Z_0)
\]
は$A/(G_1)$の三次元単純表現を与える。
その中心指標は旧論文VIの中心生成元$U,V,W$で
\[
 (g,U,V,W)=(0,-1,-4\omega^{2k},4\omega^k)
\]
であり、同じ中心点で二つの非同型類がある。
\end{proposition}
\begin{proof}
$\omega^2+\omega+1=0$上で
$[X_0,Y_0]=Z_0^2$、
$[Y_0,Z_0]=X_0^2+I$、
$[Z_0,X_0]=Y_0^2$および
$\Omega_1(X_0,Y_0,Z_0)+X_0=0$を直接確認できる。
これは付属の$\CC[\omega]/(\omega^2+\omega+1)$上の
有理数記号計算で再現する。
$Y_0$は相異なる三固有値を持ち、
$X_0$の非対角成分はすべて非零なので
共通不変部分空間は存在しない。
$g$は符号変換で反転し、$U,V,W$の規格化された
中心多項式の代入により表示した中心指標を得る。
$V,W$の非可換語による明示式を
\texttt{computations/verify\_central\_quotient\_reps.py}
に追加した。これらの中心指標は、
$\mathbb Q[\omega]/(\omega^2+\omega+1)$上の
独立した記号行列計算でも検証した。
なお保存版の全スクリプトの一括実行は未実施である。
異なる$k$は中心指標で、同じ$k$の符号は
$\operatorname{Tr}X=-\varepsilon$で区別できる。
\end{proof}

\begin{lemma}[六次最小多項式の回復]
\label{lem:local-z-minpoly}
$C=Z(R_1)$は系\ref{cor:quotient-normal-center}により正規整域であり、
中心元の定義は$U=Z^3(g-Z^3)$である。
このとき$K=\operatorname{Frac}C$上の
$Z$の最小多項式および被約特性多項式は
\[
 P(T)=T^6-gT^3+U
\]
である。
\end{lemma}
\begin{proof}
$P(Z)=0$なので最小多項式$m(T)$は$P(T)$を割る。
$m$の係数は$P$の根の対称式として$C$上整であり
$K$に属するため、$C$の正規性から$m(T)\in C[T]$。
軸上の二つの三次元表現は$g=0,U=1/7$という
同一の中心指標を持つ。ここで$U=-Z^6=1/7$は
軸上の$Z^3=-s/7$と$s^2=-7$から直接得られる。
二表現の$Z$固有値を併せると六つが相異なるので、
$m$をこの中心指標に特殊化すると六根を持つ。
従って$\deg m\ge6$であり、$m=P$。
次数六の中心的斜体での被約特性多項式は
次数六のモニック多項式で$m$の倍数だから同じものとなる。
\end{proof}

\begin{proposition}[八つの三次元単純表現の一次拡大]
\label{prop:quotient-full-ext1}
各四つの中心点の二つの三次元単純表現を$S_+,S_-$とする。
命題\ref{prop:local-model}のPBW・既知Sklyanin入力の下で
\[
 \dim\operatorname{Ext}^1_{R_1}(S_\epsilon,S_\eta)
 =\dim\operatorname{Ext}^1_{R_1/(g)}(S_\epsilon,S_\eta)
 =\begin{cases}0&\epsilon=\eta,\\2&\epsilon\ne\eta.\end{cases}
\]
異なる中心点に属する単純表現の間では、全てのExt群が零である。
従って一次拡大箙は、二頂点と各方向二矢印を持つ四つの非連結成分からなる。
\end{proposition}
\begin{proof}
命題\ref{prop:local-model}の完備化でMorita同値を取り、
$E=\operatorname{End}_R(R\oplus I)$、
$R=\CC[[p,q,a,b]]/(pq-ab)$、$I=(p,a)$とする。
ここで$\gamma=(b-a)/2$である。
$R$の極大イデアルを$\mathfrak m$とすると
\[
 I^*=p^{-1}(p,b),\qquad II^*=(p,a,b,q)=\mathfrak m.
\]
$R$と$I$は非同型直既約で、両者の自己準同型環は$R$。
従ってJacobson radicalは
\[
 J_E=\begin{pmatrix}\mathfrak m&I^*\\I&\mathfrak m\end{pmatrix},\qquad
 J_E/J_E^2\simeq
 \begin{pmatrix}0&I^*/\mathfrak m I^*\\I/\mathfrak m I&0\end{pmatrix}.
\]
$I$と$I^*$はそれぞれ二つの極小生成元を持つので、
二つの非対角成分はともに$\CC^2$である。
projective coverのradicalを使うExt$^1$の計算から、
自己拡大は零、相互拡大は二次元となる。

さらに$\gamma E\subset J_E^2$なので、$\gamma$で割っても
$J_E/J_E^2$は変化しない。$(g)=(\gamma)$であるから
中心切断$R_1/(g)$でも同じ一次拡大表を得る。
中心上有限な代数の有限長加群間のExtは同じ中心閉点に支持され、
忠実平坦完備化と可換なので、元の代数でのExtも同じである。
異なる中心指標を持つ二加群では、ある中心元が異なるスカラーで作用する。
Ext上の二つの中心作用は等しいため、その差の非零スカラーがExtを殺し、
従って全Ext群が零となる。

独立の有限計算でも同じ結果を確認した。
上三角block表現の非対角三行列を27変数として関係式を線型化した行列を$J$、
$g=Y^3-ZYX+X$の線型化を追加した行列を$J_g$とすると、
軸上・非軸上の両方について
\[
\begin{array}{c|c|c|c|c}
 &\operatorname{rank}J&\operatorname{rank}J_g
 &\operatorname{rank}(\text{内的微分})&\dim\operatorname{Ext}^1\\\hline
 \text{自己}&19&19&8&0\\
 \text{相互}&16&16&9&2
\end{array}
\]
となる。これは$\mathbb Q[t]/(q)$と
$\mathbb Q[\omega]/(\omega^2+\omega+1)$上の厳密な行列階数である。
再現コード\texttt{computations/verify\_central\_quotient\_extensions.py}は
$J$と$J_g$が内的微分を殺すことも検査する。
符号反転は各点の二表現を交換し、三乗根の対称性は非軸上三点を巡回するので、
この計算は八つの全単純表現を扱っている。
\end{proof}

\begin{proposition}[軸上の三次元表現の無限小剛性と整数小行列式]
\label{prop:quotient-axial-rigidity}
命題\ref{prop:local-two-axial}の二つの三次元単純表現$V$は、
$R_1$-加群としても$R_1/(G_1)$-加群としても
\[
 \operatorname{Ext}^1(V,V)=0
\]
を満たす。これは有限体上の非零$19$次小行列式により
特性零で証明できる。
\end{proposition}
\begin{proof}
$R_1$の三つの関係式を三行列$X,Y,Z$に代入した
$27$個の多項式を、各行列の$9$成分からなる$27$変数で
微分し、$27\times27$行列$J$を得る。
列は$X,Y,Z$それぞれに行優先で標準行列単位$E_{ij}$を並べ、
行は関係式
$[X,Y]-Z^2,\ [Y,Z]-X^2-I,\ [Z,X]-Y^2$
をそれぞれ行優先で並べる。

$q(t)=7t^6+21t^4+14t^2+1$は$\mathbb F_3$上既約であり、
従って$\mathbb Q$上既約である。
$\mathbb F_{29}$において$t=2$は$q$の根で、
\[
 (\lambda_0,\lambda_1,\lambda_2)=(3,10,2),\qquad
 (d_0,d_1,d_2)=(7,21,1)\pmod{29}.
\]
$J$の次の零始まり行番号・列番号からなる小行列を取る：
\[
\begin{split}
 I={}&(2,0,1,5,3,4,10,11,9,14,6,7,13,20,16,17,18,19,8),\\
 K={}&(0,1,2,3,4,5,6,7,8,9,10,11,12,13,14,16,18,19,20).
\end{split}
\]
整数係数の明示的なJacobi微分から計算すると
\[
 \det J_{I,K}\equiv15\not\equiv0\pmod{29}.
\]
再現コードは
\texttt{computations/verify\_central\_quotient\_reps.py}に
あり、各微分列を生成して法$29$のGaussian消去を行う。
従って$\mathbb Q[t]/(q)$上でrankは少なくとも$19$。
既約性により、すべての複素根の表現で同じ下界が成立する。

各$V$は単純なので共役軌道の次元は
$\dim\operatorname{PGL}_3=8$である。
軌道方向はJacobianの核に含まれるからrankは高々$19$。
従って接空間の次元は$27-19=8$で、軌道方向と一致する。
この接空間を内的導分で割ったものが
$\operatorname{Ext}^1_{R_1}(V,V)$なので零となる。
$G_1=0$を追加した表現スキームの接空間は前者の部分空間で
同じ共役軌道を含むため、こちらも次元は$8$であり、
$\operatorname{Ext}^1_{R_1/(G_1)}(V,V)=0$となる。
\end{proof}
